\IfFileExists{stacks-project.cls}{%
\documentclass{stacks-project}
}{%
\documentclass{amsart}
}

% For dealing with references we use the comment environment
\usepackage{verbatim}
\newenvironment{reference}{\comment}{\endcomment}
%\newenvironment{reference}{}{}
\newenvironment{slogan}{\comment}{\endcomment}
\newenvironment{history}{\comment}{\endcomment}

% For commutative diagrams we use Xy-pic
\usepackage[all]{xy}

% We use 2cell for 2-commutative diagrams.
\xyoption{2cell}
\UseAllTwocells

% We use multicol for the list of chapters between chapters
\usepackage{multicol}

% This is generally recommended for better output
\usepackage{lmodern}
\usepackage[T1]{fontenc}

% For cross-file-references
\usepackage{xr-hyper}

% Package for hypertext links:
\usepackage{hyperref}

% For any local file, say "hello.tex" you want to link to please
% use \externaldocument[hello-]{hello}
\externaldocument[introduction-]{introduction}
\externaldocument[conventions-]{conventions}
\externaldocument[sets-]{sets}
\externaldocument[categories-]{categories}
\externaldocument[topology-]{topology}
\externaldocument[sheaves-]{sheaves}
\externaldocument[sites-]{sites}
\externaldocument[stacks-]{stacks}
\externaldocument[fields-]{fields}
\externaldocument[algebra-]{algebra}
\externaldocument[brauer-]{brauer}
\externaldocument[homology-]{homology}
\externaldocument[derived-]{derived}
\externaldocument[simplicial-]{simplicial}
\externaldocument[more-algebra-]{more-algebra}
\externaldocument[smoothing-]{smoothing}
\externaldocument[modules-]{modules}
\externaldocument[sites-modules-]{sites-modules}
\externaldocument[injectives-]{injectives}
\externaldocument[cohomology-]{cohomology}
\externaldocument[sites-cohomology-]{sites-cohomology}
\externaldocument[dga-]{dga}
\externaldocument[dpa-]{dpa}
\externaldocument[sdga-]{sdga}
\externaldocument[hypercovering-]{hypercovering}
\externaldocument[schemes-]{schemes}
\externaldocument[constructions-]{constructions}
\externaldocument[properties-]{properties}
\externaldocument[morphisms-]{morphisms}
\externaldocument[coherent-]{coherent}
\externaldocument[divisors-]{divisors}
\externaldocument[limits-]{limits}
\externaldocument[varieties-]{varieties}
\externaldocument[topologies-]{topologies}
\externaldocument[descent-]{descent}
\externaldocument[perfect-]{perfect}
\externaldocument[more-morphisms-]{more-morphisms}
\externaldocument[flat-]{flat}
\externaldocument[groupoids-]{groupoids}
\externaldocument[more-groupoids-]{more-groupoids}
\externaldocument[etale-]{etale}
\externaldocument[chow-]{chow}
\externaldocument[intersection-]{intersection}
\externaldocument[pic-]{pic}
\externaldocument[weil-]{weil}
\externaldocument[adequate-]{adequate}
\externaldocument[dualizing-]{dualizing}
\externaldocument[duality-]{duality}
\externaldocument[discriminant-]{discriminant}
\externaldocument[derham-]{derham}
\externaldocument[local-cohomology-]{local-cohomology}
\externaldocument[algebraization-]{algebraization}
\externaldocument[curves-]{curves}
\externaldocument[resolve-]{resolve}
\externaldocument[models-]{models}
\externaldocument[functors-]{functors}
\externaldocument[equiv-]{equiv}
\externaldocument[pione-]{pione}
\externaldocument[etale-cohomology-]{etale-cohomology}
\externaldocument[proetale-]{proetale}
\externaldocument[relative-cycles-]{relative-cycles}
\externaldocument[more-etale-]{more-etale}
\externaldocument[trace-]{trace}
\externaldocument[crystalline-]{crystalline}
\externaldocument[spaces-]{spaces}
\externaldocument[spaces-properties-]{spaces-properties}
\externaldocument[spaces-morphisms-]{spaces-morphisms}
\externaldocument[decent-spaces-]{decent-spaces}
\externaldocument[spaces-cohomology-]{spaces-cohomology}
\externaldocument[spaces-limits-]{spaces-limits}
\externaldocument[spaces-divisors-]{spaces-divisors}
\externaldocument[spaces-over-fields-]{spaces-over-fields}
\externaldocument[spaces-topologies-]{spaces-topologies}
\externaldocument[spaces-descent-]{spaces-descent}
\externaldocument[spaces-perfect-]{spaces-perfect}
\externaldocument[spaces-more-morphisms-]{spaces-more-morphisms}
\externaldocument[spaces-flat-]{spaces-flat}
\externaldocument[spaces-groupoids-]{spaces-groupoids}
\externaldocument[spaces-more-groupoids-]{spaces-more-groupoids}
\externaldocument[bootstrap-]{bootstrap}
\externaldocument[spaces-pushouts-]{spaces-pushouts}
\externaldocument[spaces-chow-]{spaces-chow}
\externaldocument[groupoids-quotients-]{groupoids-quotients}
\externaldocument[spaces-more-cohomology-]{spaces-more-cohomology}
\externaldocument[spaces-simplicial-]{spaces-simplicial}
\externaldocument[spaces-duality-]{spaces-duality}
\externaldocument[formal-spaces-]{formal-spaces}
\externaldocument[restricted-]{restricted}
\externaldocument[spaces-resolve-]{spaces-resolve}
\externaldocument[formal-defos-]{formal-defos}
\externaldocument[defos-]{defos}
\externaldocument[cotangent-]{cotangent}
\externaldocument[examples-defos-]{examples-defos}
\externaldocument[algebraic-]{algebraic}
\externaldocument[examples-stacks-]{examples-stacks}
\externaldocument[stacks-sheaves-]{stacks-sheaves}
\externaldocument[criteria-]{criteria}
\externaldocument[artin-]{artin}
\externaldocument[quot-]{quot}
\externaldocument[stacks-properties-]{stacks-properties}
\externaldocument[stacks-morphisms-]{stacks-morphisms}
\externaldocument[stacks-limits-]{stacks-limits}
\externaldocument[stacks-cohomology-]{stacks-cohomology}
\externaldocument[stacks-perfect-]{stacks-perfect}
\externaldocument[stacks-introduction-]{stacks-introduction}
\externaldocument[stacks-more-morphisms-]{stacks-more-morphisms}
\externaldocument[stacks-geometry-]{stacks-geometry}
\externaldocument[moduli-]{moduli}
\externaldocument[moduli-curves-]{moduli-curves}
\externaldocument[examples-]{examples}
\externaldocument[exercises-]{exercises}
\externaldocument[guide-]{guide}
\externaldocument[desirables-]{desirables}
\externaldocument[coding-]{coding}
\externaldocument[obsolete-]{obsolete}
\externaldocument[fdl-]{fdl}
\externaldocument[index-]{index}

% Theorem environments.
%
\theoremstyle{plain}
\newtheorem{theorem}[subsection]{Theorem}
\newtheorem{proposition}[subsection]{Proposition}
\newtheorem{lemma}[subsection]{Lemma}

\theoremstyle{definition}
\newtheorem{definition}[subsection]{Definition}
\newtheorem{example}[subsection]{Example}
\newtheorem{exercise}[subsection]{Exercise}
\newtheorem{situation}[subsection]{Situation}

\theoremstyle{remark}
\newtheorem{remark}[subsection]{Remark}
\newtheorem{remarks}[subsection]{Remarks}

\numberwithin{equation}{subsection}

% Macros
%
\def\lim{\mathop{\mathrm{lim}}\nolimits}
\def\colim{\mathop{\mathrm{colim}}\nolimits}
\def\Spec{\mathop{\mathrm{Spec}}}
\def\Hom{\mathop{\mathrm{Hom}}\nolimits}
\def\Ext{\mathop{\mathrm{Ext}}\nolimits}
\def\SheafHom{\mathop{\mathcal{H}\!\mathit{om}}\nolimits}
\def\SheafExt{\mathop{\mathcal{E}\!\mathit{xt}}\nolimits}
\def\Sch{\mathit{Sch}}
\def\Mor{\mathop{\mathrm{Mor}}\nolimits}
\def\Ob{\mathop{\mathrm{Ob}}\nolimits}
\def\Sh{\mathop{\mathit{Sh}}\nolimits}
\def\NL{\mathop{N\!L}\nolimits}
\def\CH{\mathop{\mathrm{CH}}\nolimits}
\def\proetale{{pro\text{-}\acute{e}tale}}
\def\etale{{\acute{e}tale}}
\def\QCoh{\mathit{QCoh}}
\def\Ker{\mathop{\mathrm{Ker}}}
\def\Im{\mathop{\mathrm{Im}}}
\def\Coker{\mathop{\mathrm{Coker}}}
\def\Coim{\mathop{\mathrm{Coim}}}

% Boxtimes
%
\DeclareMathSymbol{\boxtimes}{\mathbin}{AMSa}{"02}

%
% Macros for moduli stacks/spaces
%
\def\QCohstack{\mathcal{QC}\!\mathit{oh}}
\def\Cohstack{\mathcal{C}\!\mathit{oh}}
\def\Spacesstack{\mathcal{S}\!\mathit{paces}}
\def\Quotfunctor{\mathrm{Quot}}
\def\Hilbfunctor{\mathrm{Hilb}}
\def\Curvesstack{\mathcal{C}\!\mathit{urves}}
\def\Polarizedstack{\mathcal{P}\!\mathit{olarized}}
\def\Complexesstack{\mathcal{C}\!\mathit{omplexes}}
% \Pic is the operator that assigns to X its picard group, usage \Pic(X)
% \Picardstack_{X/B} denotes the Picard stack of X over B
% \Picardfunctor_{X/B} denotes the Picard functor of X over B
\def\Pic{\mathop{\mathrm{Pic}}\nolimits}
\def\Picardstack{\mathcal{P}\!\mathit{ic}}
\def\Picardfunctor{\mathrm{Pic}}
\def\Deformationcategory{\mathcal{D}\!\mathit{ef}}
\usepackage{cleveref}
\crefname{originallemma}{original lemma}{original lemmas}
\Crefname{originallemma}{Original lemma}{Original lemmas}
\makeatletter
\newlabel{derived-lemma-two-ss-complex-functor}{{015J}{source}{}{sourcepointer-1}{}}
\newlabel{derived-lemma-two-ss-complex-functor@cref}{{[originallemma][0][]015J}{[1][0][]source}{}{}{}}
\makeatother
\crefname{originallemma}{original lemma}{original lemmas}
\Crefname{originallemma}{Original lemma}{Original lemmas}
\makeatletter
\newlabel{coherent-lemma-quasi-coherence-higher-direct-images}{{01XJ}{source}{}{sourcepointer-2}{}}
\newlabel{coherent-lemma-quasi-coherence-higher-direct-images@cref}{{[originallemma][0][]01XJ}{[1][0][]source}{}{}{}}
\makeatother
\crefname{originallemma}{original lemma}{original lemmas}
\Crefname{originallemma}{Original lemma}{Original lemmas}
\makeatletter
\newlabel{coherent-lemma-induction-principle}{{08DR}{source}{}{sourcepointer-3}{}}
\newlabel{coherent-lemma-induction-principle@cref}{{[originallemma][0][]08DR}{[1][0][]source}{}{}{}}
\makeatother
\crefname{originallemma}{original lemma}{original lemmas}
\Crefname{originallemma}{Original lemma}{Original lemmas}
\makeatletter
\newlabel{cohomology-lemma-unbounded-relative-mayer-vietoris}{{08HZ}{source}{}{sourcepointer-4}{}}
\newlabel{cohomology-lemma-unbounded-relative-mayer-vietoris@cref}{{[originallemma][0][]08HZ}{[1][0][]source}{}{}{}}
\makeatother
\crefname{originallemma}{original lemma}{original lemmas}
\Crefname{originallemma}{Original lemma}{Original lemmas}
\makeatletter
\newlabel{cohomology-lemma-cone-tor-amplitude}{{08CJ}{source}{}{sourcepointer-5}{}}
\newlabel{cohomology-lemma-cone-tor-amplitude@cref}{{[originallemma][0][]08CJ}{[1][0][]source}{}{}{}}
\makeatother
\crefname{originallemma}{original lemma}{original lemmas}
\Crefname{originallemma}{Original lemma}{Original lemmas}
\makeatletter
\newlabel{derived-lemma-leray-acyclicity}{{015E}{source}{}{sourcepointer-6}{}}
\newlabel{derived-lemma-leray-acyclicity@cref}{{[originallemma][0][]015E}{[1][0][]source}{}{}{}}
\makeatother
\crefname{originallemma}{original lemma}{original lemmas}
\Crefname{originallemma}{Original lemma}{Original lemmas}
\makeatletter
\newlabel{coherent-lemma-relative-affine-vanishing}{{01XC}{source}{}{sourcepointer-7}{}}
\newlabel{coherent-lemma-relative-affine-vanishing@cref}{{[originallemma][0][]01XC}{[1][0][]source}{}{}{}}
\makeatother
\crefname{originalequation}{original equation}{original equations}
\Crefname{originalequation}{Original equation}{Original equations}
\makeatletter
\newlabel{cohomology-equation-projection-formula-map}{{0B53}{source}{}{sourcepointer-8}{}}
\newlabel{cohomology-equation-projection-formula-map@cref}{{[originalequation][0][]0B53}{[1][0][]source}{}{}{}}
\makeatother
\crefname{originallemma}{original lemma}{original lemmas}
\Crefname{originallemma}{Original lemma}{Original lemmas}
\makeatletter
\newlabel{cohomology-lemma-variant-derived-pullback}{{08DE}{source}{}{sourcepointer-9}{}}
\newlabel{cohomology-lemma-variant-derived-pullback@cref}{{[originallemma][0][]08DE}{[1][0][]source}{}{}{}}
\makeatother
\crefname{originallemma}{original lemma}{original lemmas}
\Crefname{originallemma}{Original lemma}{Original lemmas}
\makeatletter
\newlabel{lemma-quasi-coherence-pushforward-direct-sums}{{08DZ}{source}{}{sourcepointer-10}{}}
\newlabel{lemma-quasi-coherence-pushforward-direct-sums@cref}{{[originallemma][0][]08DZ}{[1][0][]source}{}{}{}}
\makeatother
\crefname{originallemma}{original lemma}{original lemmas}
\Crefname{originallemma}{Original lemma}{Original lemmas}
\makeatletter
\newlabel{lemma-cohomology-base-change}{{08EU}{source}{}{sourcepointer-11}{}}
\newlabel{lemma-cohomology-base-change@cref}{{[originallemma][0][]08EU}{[1][0][]source}{}{}{}}
\makeatother
\crefname{originallemma}{original lemma}{original lemmas}
\Crefname{originallemma}{Original lemma}{Original lemmas}
\makeatletter
\newlabel{homology-lemma-biregular-ss-converges}{{012W}{source}{}{sourcepointer-12}{}}
\newlabel{homology-lemma-biregular-ss-converges@cref}{{[originallemma][0][]012W}{[1][0][]source}{}{}{}}
\makeatother
\title{Locally bounded flat complexes of modules over the structure sheaf have bounded, quasi-coherent, finite tor dimension pushforwards that commute with base change}
\author{The Stacks Project Authors}
\date{stacks-project:perfect.tex:local-snapshot-20261012; GNU FDL 1.2 or later}
\begin{document}
\maketitle
\setcounter{subsection}{1}
\noindent\textit{Editorial note.} All delivered mathematical text is the author's own, copied from the pinned Stacks Project chapter file perfect.tex; no statement or proof step was rewritten, reordered or removed. The counted span is the author environment and its own proof environment. References inside the span are handed to the assembler as navigation stubs that carry the canonical tag of the cited tag only; whether each of them is load-bearing is recorded on the stub itself.\par\medskip
\noindent\textit{Editorial note.} One sentence inside the counted proof reads, verbatim: 'Some details omitted.' (at the end of the argument for part (5), where the author says the induction on the length of the bounded complex finishes after the distinguished triangle is displayed). The counted span is otherwise the author's own proof of all five parts, and the omitted clause concerns the induction whose displayed triangle and base case are given in the same span; per the project rule the phrase is recorded here verbatim.\par\medskip
\noindent\textit{Allowed standard prerequisites (editorial).}\begin{itemize}
\item The first spectral sequence of a complex of functors converges and relates the higher direct images of the terms to those of the complex. invoked at proof lines 5797-5801. Produces the local boundedness and quasi-coherence of the pushforward of the complex.
\item Higher direct images of quasi-coherent modules under a quasi-compact quasi-separated morphism are quasi-coherent, and the biregular spectral sequence converges. invoked at proof lines 5799-5802. Supplies the quasi-coherence statement needed for parts (1) and (2).
\item The induction principle on quasi-compact opens reduces a statement about all quasi-compact opens to the affine case, and relative Mayer-Vietoris plus the behaviour of tor amplitude in triangles handles the union step. invoked at proof lines 5803-5815. Reduces part (3) to the affine case, where the pushforward is computed by an acyclic resolution.
\item For an affine morphism the higher direct images of quasi-coherent modules vanish, so the pushforward of the complex is the complex of the pushforwards. invoked at proof lines 5817-5820. Turns the pushforward in the affine case into a bounded complex of flat modules, giving finite tor dimension.
\item The projection formula supplies a functorial comparison map between the two sides of the base change statements. invoked at proof lines 5827-5830. Reduces parts (4) and (5) to showing that the comparison map is an isomorphism.
\item The pushforward of a quasi-coherent module commutes with direct sums, and the derived pullback of a quasi-coherent complex is computed by the tensor product. invoked at proof lines 5835-5845. Lets the proof check the commutation with direct sums that the cone argument needs for the single module case.
\end{itemize}\medskip
\ifdefined\hypertarget\hypertarget{sourcepointer-1}{}\fi
\noindent\textit{Source reference.} Original lemma 015J: sources/stacks-chow/tags.txt, source lines 1--1. Cited at proof lines 5797-5799 for the spectral sequence used to prove parts (1) and (2); the conclusion drawn from it is inside the counted span.\par\medskip
\ifdefined\hypertarget\hypertarget{sourcepointer-2}{}\fi
\noindent\textit{Source reference.} Original lemma 01XJ: sources/stacks-chow/tags.txt, source lines 1--1. Cited at proof lines 5799-5801 for the quasi-coherence input of part (2); the spectral sequence argument that consumes it is inside the counted span.\par\medskip
\ifdefined\hypertarget\hypertarget{sourcepointer-3}{}\fi
\noindent\textit{Source reference.} Original lemma 08DR: sources/stacks-chow/tags.txt, source lines 1--1. Invoked at proof lines 5803-5805 to reduce part (3) to the affine case; the union step and the affine computation are inside the counted span.\par\medskip
\ifdefined\hypertarget\hypertarget{sourcepointer-4}{}\fi
\noindent\textit{Source reference.} Original lemma 08HZ: sources/stacks-chow/tags.txt, source lines 1--1. Cited at proof lines 5807-5810 for the displayed triangle used in the union step; the tor amplitude conclusion is inside the counted span.\par\medskip
\ifdefined\hypertarget\hypertarget{sourcepointer-5}{}\fi
\noindent\textit{Source reference.} Original lemma 08CJ: sources/stacks-chow/tags.txt, source lines 1--1. Cited at proof lines 5811-5813 to pass from the pieces of the triangle to the union; the application is inside the counted span.\par\medskip
\ifdefined\hypertarget\hypertarget{sourcepointer-6}{}\fi
\noindent\textit{Source reference.} Original lemma 015E: sources/stacks-chow/tags.txt, source lines 1--1. Cited at proof lines 5817-5819 for the identification of the pushforward of the complex with the complex of pushforwards in the affine case; the flatness discussion that follows is inside the counted span.\par\medskip
\ifdefined\hypertarget\hypertarget{sourcepointer-7}{}\fi
\noindent\textit{Source reference.} Original lemma 01XC: sources/stacks-chow/tags.txt, source lines 1--1. Cited with the previous label at proof lines 5818-5820 for the same identification; the finiteness of tor dimension concluded from it is inside the counted span.\par\medskip
\ifdefined\hypertarget\hypertarget{sourcepointer-8}{}\fi
\noindent\textit{Source reference.} Original equation 0B53: sources/stacks-chow/tags.txt, source lines 1--1. Cited at proof lines 5827-5830 to obtain the comparison map whose bijectivity is the content of part (5); the reduction to the affine case and the direct sum argument are inside the counted span.\par\medskip
\ifdefined\hypertarget\hypertarget{sourcepointer-9}{}\fi
\noindent\textit{Source reference.} Original lemma 08DE: sources/stacks-chow/tags.txt, source lines 1--1. Cited at proof lines 5838-5841 to identify the tensor product in the single degree case; the commutation with direct sums drawn from it is inside the counted span.\par\medskip
\ifdefined\hypertarget\hypertarget{sourcepointer-10}{}\fi
\noindent\textit{Source reference.} Original lemma 08DZ: sources/stacks-perfect/pinned-original-source.tex, source lines 803--803. Used at proof lines 5843-5845 for the single degree case of the direct sum argument; the induction that extends it to the whole complex is inside the counted span.\par\medskip
\ifdefined\hypertarget\hypertarget{sourcepointer-11}{}\fi
\noindent\textit{Source reference.} Original lemma 08EU: sources/stacks-perfect/pinned-original-source.tex, source lines 5269--5269. The proof states at proof lines 5831-5834 that the rest is the same as the proof of that statement except for the direct sum argument; the direct sum argument specific to this lemma is inside the counted span.\par\medskip
\ifdefined\hypertarget\hypertarget{sourcepointer-12}{}\fi
\noindent\textit{Source reference.} Original lemma 012W: sources/stacks-chow/tags.txt, source lines 1--1. Cited at proof lines 5799-5802 together with the quasi-coherence statement to conclude parts (1) and (2); the spectral sequence argument that consumes the convergence is inside the counted span.\par\medskip
\begin{lemma}
\label{lemma-cohomology-de-rham-base-change}
Let $a : X \to S$ be a quasi-compact and quasi-separated morphism
of schemes. Let $\mathcal{F}^\bullet$ be a locally bounded
complex of $a^{-1}\mathcal{O}_S$-modules. Assume for all $n \in \mathbf{Z}$
the sheaf $\mathcal{F}^n$ is a flat $a^{-1}\mathcal{O}_S$-module and
$\mathcal{F}^n$ has the structure of a quasi-coherent $\mathcal{O}_X$-module
compatible with the given $a^{-1}\mathcal{O}_S$-module structure (but the
differentials in the complex $\mathcal{F}^\bullet$ need not
be $\mathcal{O}_X$-linear). Then the following hold
\begin{enumerate}
\item $Ra_*\mathcal{F}^\bullet$ is locally bounded,
\item $Ra_*\mathcal{F}^\bullet$ is in $D_\QCoh(\mathcal{O}_S)$,
\item $Ra_*\mathcal{F}^\bullet$ locally has finite tor dimension,
\item $\mathcal{G} \otimes_{\mathcal{O}_S}^\mathbf{L} Ra_*\mathcal{F}^\bullet =
Ra_*(a^{-1}\mathcal{G} \otimes_{a^{-1}\mathcal{O}_S} \mathcal{F}^\bullet)$
for $\mathcal{G} \in \QCoh(\mathcal{O}_S)$, and
\item $K \otimes_{\mathcal{O}_S}^\mathbf{L} Ra_*\mathcal{F}^\bullet =
Ra_*(a^{-1}K \otimes_{a^{-1}\mathcal{O}_S}^\mathbf{L} \mathcal{F}^\bullet)$
for $K \in D_\QCoh(\mathcal{O}_S)$.
\end{enumerate}
\end{lemma}

\begin{proof}
Parts (1), (2), (3) are local on $S$ hence we may and do assume $S$
is affine. Since $a$ is quasi-compact we conclude that $X$ is quasi-compact.
Since $\mathcal{F}^\bullet$ is locally bounded, we conclude that
$\mathcal{F}^\bullet$ is bounded.

\medskip\noindent
For (1) and (2) we can use the first spectral sequence
$R^pa_*\mathcal{F}^q \Rightarrow R^{p + q}a_*\mathcal{F}^\bullet$ of
Derived Categories, Lemma \ref{derived-lemma-two-ss-complex-functor}.
Combining Cohomology of Schemes, Lemma
\ref{coherent-lemma-quasi-coherence-higher-direct-images}
and Homology, Lemma \ref{homology-lemma-biregular-ss-converges}
we conclude.

\medskip\noindent
Let us prove (3) by the induction principle of
Cohomology of Schemes, Lemma \ref{coherent-lemma-induction-principle}.
Namely, for a quasi-compact open of $U$ of $X$ consider the
condition that $R(a|_U)_*(\mathcal{F}^\bullet|_U)$ has
finite tor dimension. If $U, V$ are quasi-compact open in
$X$, then we have a relative Mayer-Vietoris distinguished triangle
$$
R(a|_{U \cup V})_*\mathcal{F}^\bullet|_{U \cup V} \to
R(a|_U)_*\mathcal{F}^\bullet|_U \oplus
R(a|_V)_*\mathcal{F}^\bullet|_V \to
R(a|_{U \cap V})_*\mathcal{F}^\bullet|_{U \cap V} \to
$$
by Cohomology, Lemma \ref{cohomology-lemma-unbounded-relative-mayer-vietoris}.
By the behaviour of tor amplitude in distinguished triangles
(see Cohomology, Lemma \ref{cohomology-lemma-cone-tor-amplitude})
we see that if we know the result for $U$, $V$, $U \cap V$, then
the result holds for $U \cup V$. This reduces us to the case where
$X$ is affine. In this case we have
$$
Ra_*\mathcal{F}^\bullet = a_*\mathcal{F}^\bullet
$$
by Leray's acyclicity lemma
(Derived Categories, Lemma \ref{derived-lemma-leray-acyclicity})
and the vanishing of higher direct images of quasi-coherent modules
under an affine morphism
(Cohomology of Schemes, Lemma \ref{coherent-lemma-relative-affine-vanishing}).
Since $\mathcal{F}^n$ is $S$-flat by assumption and $X$ affine, the modules
$a_*\mathcal{F}^n$ are flat for all $n$. Hence $a_*\mathcal{F}^\bullet$
is a bounded complex of flat $\mathcal{O}_S$-modules and hence has
finite tor dimension.

\medskip\noindent
Proof of part (5). Denote
$a' : (X, a^{-1}\mathcal{O}_S) \to (S, \mathcal{O}_S)$
the obvious flat morphism of ringed spaces. Part (5) says that
$$
K \otimes_{\mathcal{O}_S}^\mathbf{L} Ra'_*\mathcal{F}^\bullet =
Ra'_*(L(a')^*K \otimes_{a^{-1}\mathcal{O}_S}^\mathbf{L}
\mathcal{F}^\bullet)
$$
Thus
Cohomology, Equation (\ref{cohomology-equation-projection-formula-map})
gives a functorial map from the left to the right and we want to show
this map is an isomorphism.
This question is local on $S$ hence we may and do assume $S$
is affine. The rest of the proof is {\it exactly} the same as the
proof of Lemma \ref{lemma-cohomology-base-change} except that we have
to show that the functor
$K \mapsto Ra'_*(L(a')^*K \otimes_{a^{-1}\mathcal{O}_S}^\mathbf{L}
\mathcal{F}^\bullet)$ commutes with direct sums.
This is where we will use $\mathcal{F}^n$ has the structure
of a quasi-coherent $\mathcal{O}_X$-module. Namely, observe that
$K \mapsto L(a')^*K
\otimes_{a^{-1}\mathcal{O}_S}^\mathbf{L} \mathcal{F}^\bullet$
commutes with arbitrary direct sums. Next, if
$\mathcal{F}^\bullet$ consists of a single quasi-coherent
$\mathcal{O}_X$-module $\mathcal{F}^\bullet = \mathcal{F}^n[-n]$
then we have $L(a')^*G
\otimes_{a^{-1}\mathcal{O}_S}^\mathbf{L} \mathcal{F}^\bullet =
La^*K \otimes_{\mathcal{O}_X}^\mathbf{L} \mathcal{F}^n[-n]$, see
Cohomology, Lemma \ref{cohomology-lemma-variant-derived-pullback}.
Hence in this case the commutation with direct sums follows from
Lemma \ref{lemma-quasi-coherence-pushforward-direct-sums}.
Now, in general, since $S$ is affine (hence $X$ quasi-compact)
and $\mathcal{F}^\bullet$ is locally bounded, we see that
$$
\mathcal{F}^\bullet = (\mathcal{F}^a \to \ldots \to \mathcal{F}^b)
$$
is bounded. Arguing by induction on $b - a$ and considering the
distinguished triangle
$$
\mathcal{F}^b[-b] \to (\mathcal{F}^a \to \ldots \to \mathcal{F}^b)
\to (\mathcal{F}^a \to \ldots \to \mathcal{F}^{b - 1}) \to
\mathcal{F}^b[-b + 1]
$$
the proof of this part is finished. Some details omitted.

\medskip\noindent
Proof of part (4). Let $a' : (X, a^{-1}\mathcal{O}_S) \to (S, \mathcal{O}_S)$
be as above. Since $\mathcal{F}^\bullet$ is a locally bounded
complex of flat $a^{-1}\mathcal{O}_S$-modules we see the complex
$a^{-1}\mathcal{G} \otimes_{a^{-1}\mathcal{O}_S} \mathcal{F}^\bullet$
represents $L(a')^*\mathcal{G}
\otimes_{a^{-1}\mathcal{O}_S}^\mathbf{L}
\mathcal{F}^\bullet$ in $D(a^{-1}\mathcal{O}_S)$. Hence (4)
follows from (5).
\end{proof}

\end{document}
